% Standalone section fragment. Uses only amsmath/amssymb and ordinary LaTeX.
% Statement identifiers B-01 etc are stable and used in companion CSV.
\subsection{B.1. The corrected three-data theorem}

We use simplicial objects indexed by all order-preserving maps in $\Delta$,
including degeneracies. This is the convention explicitly declared in cote
136, scans 4 and 7. For a divided-power triple $(S,J,t)$, put
\[
 A_n=S\{U_0,\ldots,U_{n-1}\}\otimes_S
       \Lambda_S(dU_0,\ldots,dU_{n-1}),\qquad U_r=X_0+\cdots+X_r.
\]
The differential is the relative divided-power de Rham differential; in
particular $d(s)=0$ for $s\in S$. The two evaluations on $A_1$ are
$\alpha(X)=0$ and $\beta(X)=t$, and kill $dX$. A morphism of the
underlying differential rings means a unital additive and multiplicative
map commuting with $d$, without an assumption about exterior degree.
The scalar divided-power ideal $J$ need not be generated by $t$.

\paragraph{B-01. Proved: the three-data classification.}
Let $(S,J,t)$ and $(S',J',t')$ be arbitrary divided-power triples over
commutative rings. Morphisms $\Phi:A_\bullet\to A'_\bullet$ of underlying
simplicial differential rings are in bijection with data
\[
 u:S\longrightarrow S',\qquad A\in (X)_{\rm PD}\subset S'\{X\},
 \qquad B_i\in S'\{X\}\quad(i\geq1),
\]
where $u$ is a ring homomorphism, subject to
\begin{align*}
 u(t^{(i)})&=u(t)^{(i)} &&(i\geq1),\\
 \beta'(A)&=u(t),\\
 2B_i L&=0 &&(i\geq1),\\
 \binom{i+j}{i}B_{i+j}&=A^{(i)}B_j+A^{(j)}B_i &&(i,j\geq1).
 \tag{B-M}
\end{align*}
Here $L\subset S'$ is the ideal generated by all coefficients of $A$
and all $B_i$. The second line ensures $u(t)\in J'$, so the first line
is well typed. The symbol $L$ avoids confusing this coefficient ideal
with the given divided-power ideal $J'$.
The map is
\[
 \Phi_n(U_r^{(i)})=A^{(i)}(U_r)+B_i(U_r)dU_r,
 \qquad \Phi_n(dU_r)=A_X'(U_r)dU_r,
 \qquad \Phi_n|_S=u.
 \tag{B-1}
\]
The ranges are $0\leq r<n$ and $i\geq1$; the apparent $U_n$ on scan
164 is an index slip. These formulas and conditions include $t=0$,
characteristics $2$ and $3$, and rings with torsion.

\emph{Proof.}
Write $\Phi_1(X^{(i)})=A_i+B_i dX$. Compatibility with $\alpha$ gives
$A_i(0)=0$. Compatibility with $d$ gives
$\partial_X A_i=A_{i-1}\partial_X A_1$, with $A_0=1$.
The map $\partial_X:(X)_{\rm PD}\to S'\{X\}$ is an isomorphism of
$S'$-modules, since it sends $X^{(j)}$ to $X^{(j-1)}$. Induction and
$d\gamma_i(A)=\gamma_{i-1}(A)dA$ therefore give $A_i=A_1^{(i)}$.
Set $A=A_1$. Applying $\beta'$ gives the first two conditions.
The coefficient of $dX$ in
$\Phi_1(X^{(i)}X^{(j)})=\Phi_1(X^{(i)})\Phi_1(X^{(j)})$
is exactly (B-M).

In dimension two use independent variables $U,V$. The source elements
$U^{(i)}$ and $V^{(j)}$ commute. Their images have commutator
$2B_i(U)B_j(V)dU\,dV$, up to the harmless convention for the order of
the commutator. The source elements $U^{(i)}$ and $dV$ commute; the
corresponding commutator is $2B_i(U)A_X'(V)dU\,dV$.
Independence of the divided-power monomial bases in $U,V$ says precisely
that every coefficient of $2B_i$ annihilates every coefficient of $A$
and of every $B_j$. This is $2B_iL=0$. Thus no uncertain exponent in
the transcription can replace the factor $2$: the necessary equation
is established directly from the presentation.

Conversely, in one variable (B-M) verifies all divided-power
multiplication relations. The image of $dX$ has square zero and
commutes with every image of $X^{(i)}$. In several variables the
condition $2B_iL=0$ verifies the additional mixed commutation relations;
the images of different $dU_r$ anticommute and each has square zero.
Hence (B-1) defines a ring map by the presentation of $A_n$.
To check differential compatibility on every monomial, differentiate
the expanded product of the factors
$A^{(i_r)}(U_r)+B_{i_r}(U_r)dU_r$ and the selected $A_X'(U_s)dU_s$.
The one-variable equalities give the required terms. Any sign
introduced by moving a new differential past an odd summand differs
from the source sign by a multiple of
$2B_i(U_r)A_X'(U_s)$ or $2B_i(U_r)B_j(U_s)$, which vanishes by the
coefficient condition. This proves $d\Phi_n=\Phi_n d$ without
incorrectly treating $d$ as an ungraded derivation.

For a monotone map $\theta:[m]\to[n]$, each initial-segment
characteristic function $U_r$ pulls back to an initial-segment
characteristic function, $0$, or the constant $t$. Thus all
simplicial identities reduce to the displayed formulas and to the
two endpoint conditions. Conversely the $U_r^{(i)},dU_r$ are
pullbacks from the one-simplex and generate $A_n$. This proves
existence, naturality, and uniqueness.\hfill$\square$

\paragraph{B-02. Refuted, with exact replacement: omitted conditions and
unqualified divided-power preservation.}
The three conditions on scan 162, without (B-M), are false as a
classification. Over $k=\mathbf F_2$, take $S=S'=k\{T\}$, $u=1$,
$A=X$, $B_2=1$, and all other $B_i=0$. The three printed conditions
hold, whereas the relation for $(i,j)=(1,2)$ has defect $X\,dX\ne0$.
The exact replacement is B-01.

If $S'$ is torsion free over $\mathbf Z$, (B-M) implies
$B_n=A^{(n-1)}B_1$: for $i=1,j=n-1$ the induction hypothesis reduces
the equation to $n(B_n-A^{(n-1)}B_1)=0$.
This proves preservation of the canonical variable divided powers.
Preservation of the whole scalar divided-power structure requires
$u:(S,J)\to(S',J')$ to be a divided-power morphism, or follows if
$J=(t)_{\rm PD}$. It does not follow for an arbitrary larger $J$.
For example, take $S=S'=\mathbf Z[\epsilon]/(\epsilon^2)$ and
$t=t'=0$, give the source the square-zero divided-power ideal
$(\epsilon)$ and the target the zero ideal, and use the identity on
the underlying rings. The simplicial differential algebras are
identical, but the scalar ideal is not preserved.

The unqualified margin assertion on scans 170 and 172 that the
condition disappears when $2$ is invertible is also false. Take
$k=\mathbf F_3$, $S=S'=k\{T\}$, $u(T^{(n)})=0$ for $n>0$,
$A=0$, $B_3=T^{(5)}$, and all other $B_i=0$.
In (B-M) only $i+j=3$ can contribute and its two binomial coefficients
are divisible by $3$. The coefficient ideal is $(T^{(5)})$, whose
square is zero since $\binom{10}{5}=252=0$ in $k$.
All hypotheses of B-01 hold, but $X^{(3)}$ maps to the nonzero odd
form $T^{(5)}dX$. The coefficient ring $k$ is reduced; the divided-power
algebra $S'$ itself is not reduced. The strongest simple replacement
is: if $L$ contains a regular element and $2$ is invertible, all $B_i$
vanish. Every assertion in these examples is checked either by the
presentation above or by \texttt{code/B\_finite\_checks.py}.

\subsection{B.1 continued. The complete automorphism functor}

From now on completion means the separated total-weight completion,
with $T,U_r,dU_r$ all of weight one. Morphisms between completed
objects are required to be continuous for this topology. This
requirement is also automatic for natural automorphism families
over all coefficient algebras, as B-17 proves. It is false for
isolated abstract automorphisms at a fixed ring: B-15 constructs
discontinuous examples explicitly. We denote the formal
coefficient-ring functor by $\Omega_{\rm cont}$ and distinguish it
from $\Omega_{\rm abs}$ when topology and coefficient naturality
are both omitted.

\paragraph{B-03. Proved under the stated topology: exact parameters,
group law, and determinant criterion.}
Put $\widehat S=k\{\{T\}\}$ and $R=k\{\{T,X\}\}$. The elements of
$\Omega_{\rm cont}(k)$ are exactly the pairs
\[
 A=\sum_{j\geq1,h\geq0}a_{jh}X^{(j)}T^{(h)},\quad a_{10}\in k^*,
 \qquad (c_m)_{m\geq0},
 \tag{B-2}
\]
where $c_m\in {}_2R=\{c:2c=0\}$ and $c_m\to0$ in total weight.
They give B-01 data with
\[
 u(T)=A(T,T),\qquad B_{2^m}=c_m,
 \qquad B_n=\sum_{\substack{m\geq0\\\text{$m$th binary digit of $n$ is }1}}
                  A^{(n-2^m)}c_m.
 \tag{B-3}
\]
All coefficients in the last formula act on the $2$-torsion module,
so the odd multinomial coefficients in the binary divided-power
factorization are equal to one there.

Here is the full group law, also valid for composable data over
different triples whenever the indicated sums converge. Write
$\overline\Phi:R\to R'$ for the semilinear divided-power substitution
with scalar map $u$ and $X\mapsto A$, and let
$D_\Phi:R\to R'$ be its derivation, zero on scalars and satisfying
$D_\Phi(X^{(i)})=B_i$. For $\Psi=(v,C,(D_i))$, the composite
$\Phi\circ\Psi$ has data
\begin{align*}
 u_{\Phi\Psi}&=u\circ v,\\
 A_{\Phi\Psi}&=\overline\Phi(C),\\
 B_{\Phi\Psi,i}&=D_\Phi(C^{(i)})+
                  \overline\Phi(D_i)A_X'.
 \tag{B-4}
\end{align*}
The scalar map in a $k$-linear completed morphism is determined by
$A$ as in (B-3). The determinant is $\Delta=a_{10}$, and a continuous
endomorphism is an automorphism if and only if $\Delta\in k^*$.

\emph{Proof.}
For an automorphism the scalar and one-variable inverse imply that
the linear coefficient is a unit. Equivalently determinants multiply,
as is seen by taking the coefficient of $T$ in the composite scalar
substitution. Then $a_1(T)$ is a unit of $\widehat S$, so the coefficient
ideal condition of B-01 gives $2B_i=0$.
A derivation with values in a $2$-torsion module factors through
$\Gamma_{k/2}(X)$. There is the explicit presentation
\[
 \Gamma_{k/2}(X)=\bigotimes_{m\geq0}
       (k/2)[y_m]/(y_m^2),\qquad y_m=X^{(2^m)}.
\]
The binary monomial basis proves this isomorphism. The derivative
of every relation $y_m^2$ is zero, so arbitrary values $c_m$ determine
exactly one derivation. The binary product rule is (B-3).
Continuity forces $B_n\to0$, because $X^{(n)}\to0$ and
$A^{(n)}$ has weight at least $n$. Conversely $c_m\to0$ implies
$B_n\to0$: in every summand of (B-3), either $m$ is large, making
$c_m$ have arbitrarily high order, or $n-2^m$ is large.
This proves both the necessity and sufficiency of the convergence
condition and extends the polynomial classification to completions.

For the converse determinant criterion, the exterior-degree-zero
part is the simultaneous formal divided-power substitution
\[
 T\longmapsto u(T),\qquad U_r\longmapsto A(T,U_r).
\]
Its associated graded in total weight multiplies every monomial
of weight $w$ by $\Delta^w$. If $\Delta$ is a unit this is invertible.
Successively solving in weights gives a unique continuous inverse:
at weight $w$ the correction is found by applying the inverse of
multiplication by $\Delta^w$ to the current error. Completeness gives
existence and separatedness gives uniqueness. This argument uses no
division by factorials.

Remove this invertible degree-zero map from $\Phi_n$. The remaining
linear endomorphism is $1+N$ where $N$ strictly raises exterior
degree. Since there are only $n$ independent differentials,
$N^{n+1}=0$. Its inverse is the finite sum
$1-N+\cdots+(-N)^n$. Thus $\Phi_n$ is bijective. Its inverse is
a ring map, commutes with $d$, and is simplicial, by applying the
inverse to the corresponding identities for $\Phi$. It is continuous.
Finally
\[
 \Phi(f+g\,dX)=\overline\Phi(f)+
       \bigl(D_\Phi(f)+\overline\Phi(g)A_X'\bigr)dX
\]
proves (B-4), including all coefficient-semilinear terms.\hfill$\square$

\paragraph{B-04. Proved: the exact semidirect product and the
representability correction.}
Let $P(k)$ be the group of $A$ in (B-2), with all $B_i=0$, and
let $E(k)$ be the group with $A=X$, equivalently the continuous
$\widehat S$-derivations $R\to{}_2R$. Then
\[
 \Omega_{\rm cont}=E\rtimes P,
 \qquad P=P_1\rtimes\mathbf G_m,
 \tag{B-5}
\]
where $P_1$ consists of $A=X+$ terms of total weight at least two.
The group $E$ is additive, and as a functor on rings it is
\[
 E(k)=\left\{(c_{m,j,h}):2c_{m,j,h}=0,\quad
       \text{for each $(j,h)$ only finitely many $m$ occur}\right\}.
 \tag{B-6}
\]
Conjugation by $\alpha\in P$, with substitution $\overline\alpha$,
is
\[
 D\longmapsto A_X'\,\overline\alpha D\overline\alpha^{-1}.
 \tag{B-7}
\]
In particular $\lambda\in\mathbf G_m$ acts on $c_{m,j,h}$ by
$\lambda^{j+h+1-2^m}$. The successive weight quotients of $P_1$
are additive groups; thus $P_1$ is pro-unipotent in this explicit
filtered sense.

\emph{Proof.}
Taking exterior-degree zero gives a split homomorphism to $P$;
its kernel is $E$. On the one-variable generators, composition of
two kernel elements adds their derivations because $(dX)^2=0$;
the classification then proves the assertion in every simplicial
degree. Formula (B-7) follows by conjugating
$f\mapsto f+D(f)dX$. The determinant map has the section
$A=\lambda X$, proving the second decomposition. For elements tangent
to the identity, composition and inversion modulo the next weight
are addition and negation of the first nonzero homogeneous
correction. This proves the filtration assertion.

The continuity condition in (B-6) is exactly the same as
$c_m\to0$, since there are only finitely many monomials below any
fixed total weight. For every choice of a finite support set in
$m$ for each row $(j,h)$, the arrays with those support bounds are
represented by an affine scheme with coordinates satisfying
$2c_{m,j,h}=0$. Their directed union, under closed inclusions, gives
an ind-affine presentation of $E$ and hence of $\Omega_{\rm cont}$.
This functor is not affine. Indeed let $k=\prod_{m\geq0}\mathbf F_2$.
In the $m$th factor choose the kernel element with $c_m=1$ and all
other $c_r=0$. A lift of this family to $E(k)$ would have
$c_m=e_m$, the nonzero $m$th idempotent, in the constant output row
for every $m$, contradicting row finiteness. An affine scheme's
functor of points commutes with products of rings. The same example,
with $A=X$ in all factors, disproves affineness of $\Omega_{\rm cont}$.
It corrects both the unrestricted parameter list on scan 178 and its
affine-scheme conclusion.\hfill$\square$

\subsection{B.2. The complete cohomological action}

\paragraph{B-05. Proved: linear maps between the cycle objects.}
Let $A_w^q=\Gamma_k^{w-q}\Phi\otimes_k\Lambda_k^q\Psi$
be the universal weight-$w$ strand, and set
$Z_w^i=\ker(d:A_w^i\to A_w^{i+1})$, for $0\leq i\leq w$.
As a simplicial $k$-module, $Z_w^i$ is simplicially $k$-linearly
homotopy equivalent to $K(k,i)$. Consequently every simplicial
$k$-linear map $Z_w^i\to Z_v^j$ with $j\ne i$ is nullhomotopic.
The completed cycle module
$\widehat Z^i=\prod_{w\geq i}Z_w^i$ is likewise simplicially
$k$-linearly homotopy equivalent to
$K(\prod_{w\geq i}k,i)$. Thus every simplicial $k$-linear map
$\widehat Z^i\to\widehat Z^j$, $i\ne j$, is nullhomotopic,
whether or not the map is continuous.

\emph{Proof.}
The exact weight strand
\[
 0\longrightarrow kT^{(w)}\longrightarrow A_w^0
 \longrightarrow A_w^1\longrightarrow\cdots
 \longrightarrow A_w^w\longrightarrow0
 \tag{B-8}
\]
is the divided-power Poincare complex in each simplicial degree.
It is split over $\mathbf Z$ before base change to $k$: in one
variable $U^{(a)}dU\mapsto U^{(a+1)}$ is an integral contracting
operator, and tensoring the one-variable contractions gives a
contraction in any finite number of variables. It preserves total
weight. Therefore every $(Z_w^i)_n$ is a finite free $k$-module.
The normalized module in each degree is a direct summand and is
projective.

For $q<w$, the factor $\Gamma_k^{w-q}\Phi$ is simplicially
contractible. A simplicial contraction of $\Phi$ remains a
contraction after applying the positive-degree functor $\Gamma^{w-q}$;
tensor its homotopy operators with the degeneracies of
$\Lambda^q\Psi$. Thus $A_w^q$ is contractible.
The short exact sequences obtained from (B-8) show successively
that $Z_w^i$ has the single homotopy module $k$ in degree $i$.
By the Dold--Kan theorem, in the precise form of the Stacks Project,
Theorem 14.24.3 (tag 019D), its normalized complex is a bounded-below
complex of projectives with this homology.
Such a complex is chain-homotopy equivalent to $k[i]$.
For completeness, split the surjection onto degree-zero boundaries,
then proceed upwards; at the sole nonzero homology degree split off
the projective summand $k$, and continue. This writes the complex as
$k[i]$ plus elementary contractible two-term summands. Maps
$k[i]\to k[j]$ of chain complexes are zero if $i\ne j$.
Transporting the contractions through Dold--Kan proves the claim;
the equivalence of the two homotopy relations is stated explicitly
in \cite[Theorem~3.31, p.~225]{BDoldPuppe1961}.
For the completed assertion, normalization commutes with products,
since its degree-$n$ module is defined by finitely many kernels.
Take the products of the chosen chain maps and chain homotopies
for the individual weights. Their homotopy identities hold
coordinatewise, giving an actual chain-homotopy equivalence
$N\widehat Z^i\simeq(\prod_{w\geq i}k)[i]$.
No assertion that this product is projective is needed. A chain
map between these concentrated complexes in distinct degrees is
zero; conjugating by the homotopy equivalences proves the
completed assertion for arbitrary algebraic linear maps.
\hfill$\square$

\paragraph{B-06. Proved, even without continuity: the action is scalar
substitution, with no positive-degree component.}
For every commutative ring $k$ and every simplicial set $X$, there is
the weightwise comparison
\[
 H^i\operatorname{Hom}(X,\widehat A_k)
   =\prod_{w\geq i}H^i(X;k)\,T^{(w)}.
 \tag{B-9}
\]
If $\Phi\in\Omega_{\rm abs}(k)$ has scalar map $u$ and
\[
 u(T^{(w)})=\sum_{v\geq w}u_{vw}T^{(v)},
\]
then its induced operation on (B-9) is
\[
 (x_w)_{w\geq i}\longmapsto
       \left(\sum_{i\leq w\leq v}u_{vw}x_w\right)_{v\geq i}.
 \tag{B-10}
\]
Every component raising exterior/cohomological degree is zero.
The map is $u$-semilinear over $\widehat S$. Its action on the
weight-$w$ associated graded is multiplication by $\Delta^w$,
and on the lowest-weight quotient in cohomological degree $i$
it is multiplication by $\Delta^i$.

\emph{Proof.}
Products commute with kernels and images in the category of modules;
therefore taking the weightwise product gives (B-9) from the
weightwise comparison, including infinite $X$. For $i>0$, boundaries
in $\operatorname{Hom}(X,Z_w^i)$ are exactly the nullhomotopic maps.
Indeed $A_w^{i-1}\twoheadrightarrow Z_w^i$ is a degreewise
surjection of simplicial modules, hence a Kan fibration, and
$A_w^{i-1}$ is contractible. A nullhomotopy lifts along this fibration
starting with the zero lift; conversely any map through its
contractible source is nullhomotopic. For $i=0$ this is the usual
identification of maps to a constant simplicial module with
locally constant functions. Hence the cohomology comparison is also
$[X,Z_w^i]$.

For the completed cycle objects, boundaries are again precisely
nullhomotopic maps: use the surjection
$\prod_{w\geq i}A_w^{i-1}\twoheadrightarrow\widehat Z^i$,
whose source has the product of the linear contractions above.
Commutation with $d$ makes each exterior-degree component a
simplicial $k$-linear map
$\widehat Z^i\to\widehat Z^j$. For $j\ne i$ this map is
nullhomotopic by the completed assertion of B-05, so its operation
is zero for every $X$. Neither continuity of $\Phi$ nor finiteness
of $X$ is used here.

The exterior-degree-zero part $\overline\Phi$ is the continuous
substitution of B-16, even for an abstract $\Phi$, and preserves
total-weight filtration. Each matrix entry from weight $w$ to weight $v$
is a map from (B-8) to its weight-$v$ version, over the scalar map
$kT^{(w)}\to kT^{(v)}$ given by $u_{vw}$. Naturality of the short
exact-sequence connecting maps identifies its effect on cohomology
with $u_{vw}$ in every degree $i\leq w$. This is (B-10), where the
sum is finite because $u_{vw}=0$ for $v<w$.
Semilinearity follows already on cochains from
$\Phi(sa)=u(s)\Phi(a)$. The leading coefficient of
$u(T^{(w)})$ is $\Delta^w$.\hfill$\square$

\paragraph{B-07. Refuted: the literal determinant-only assertion L4.}
Take $k=\mathbf Z$ and $X=\Delta[0]$. Choose
$A=X+X^{(2)}$, all $B_i=0$. B-03 gives a continuous automorphism
with $\Delta=1$, but its effect on
$H^0(\operatorname{Hom}(X,\widehat A_k))=\mathbf Z\{\{T\}\}$
is $T\mapsto T+T^{(2)}$, which is not the identity.
Thus the literal full action does not factor through $\Delta$.
The exact replacement is B-06. The finite certificate computes the
integral unipotent substitution matrix and its inverse through
weight 12; the proof in B-03 establishes the full infinite inverse.

\subsection{B.3 and B.5. Which cohomology operations occur?}

\paragraph{B-08. Proved: the B-automorphisms induce only the identity.}
Over $\mathbf F_2$, every B-automorphism ($A=X$, $u=1$) acts as
the identity on (B-9) and on every recovered ordinary cohomology
group $H^i(X;\mathbf F_2)$. In particular none of its positive-degree
components is a nonzero Steenrod square. The same vanishing holds
over every ring, including the abstract kernel $W$ of B-16 as well
as the continuous B subgroup.

\emph{Proof.} This is B-06 with $u=1$. It also follows directly by
applying B-05 to each positive-degree component.\hfill$\square$

\paragraph{B-09. Refuted: the sphere computation proposed in RP3.}
The assertion in \texttt{study/RESEARCH\_PACKAGES.md}, RP3, that an
operation $H^i(-;\mathbf F_2)\to H^{i+1}(-;\mathbf F_2)$ is decided
on $\Delta[i+1]/\partial\Delta[i+1]$ is false for $i>0$:
the source group on that sphere is zero. For example $\mathrm{Sq}^1$
and the zero operation agree there, but disagree on a generator
$a\in H^1(\mathbf{RP}^2;\mathbf F_2)$, where
$\mathrm{Sq}^1a=a^2\ne0$. We use Steenrod--Epstein,
\emph{Cohomology Operations}, Chapter I, \S1, axiom 3, p.~1,
for the top-square identity (the source was opened).

The finite certificate supplies the ten triangles
\[
 012,013,024,035,045,125,134,145,234,235
\]
on vertices $0,\ldots,5$. It verifies that every edge has two
incident triangles and every vertex link is a five-cycle, enumerates
all $2^{15}$ one-cochains, finds $64$ one-cocycles and $32$
one-coboundaries, and produces a one-cocycle whose Alexander--Whitney
square evaluates to one on the sum of all triangles.
Thus the nonzero square is checked directly on a finite simplicial
set. The exact replacement for the faulty reduction is evaluation
on the universal cocycle of $K(\mathbf F_2,i)$; for a degree-$j$
output, its finite $(j+1)$-skeleton suffices.

\paragraph{B-10. Proved under an explicit interpretation of
$\tau_\omega$: strict natural automorphisms still yield no positive
operation.}
Let $q$ be either the Serre quotient of nonnegatively graded
$\Gamma_k(T)$-modules by bounded-weight modules, or the quotient by
the localizing Serre subcategory they generate. Let
$L_q$ be the exact reconstruction functor of Part A.2, so that
$L_q q(M\otimes_k S)=M$. Consider the functor on simplicial sets
\[
 X\longmapsto q\left(\bigoplus_w
           \operatorname{Hom}(X,A_w^\bullet)\right).
 \tag{B-11}
\]
This is the weight-graded interpretation required by the category
$\mathcal M_\omega$ on scan 182. The transcription's uncertain hat
on $\mathcal A$ makes no difference to its homogeneous pieces:
in either reading we take their graded direct sum before applying
$q$. A completion as an ungraded product is not itself an object
of this graded quotient category. The grading subsequently
forgotten on that page is the exterior grading of the differential
algebra; the ambient graded-$S$ module category remains in place.
A strict natural automorphism of its unital differential-ring
objects, with exterior degree forgotten but the graded-$S$ module
category retained, induces only the identity operation on
$H^*(X;\mathbf F_2)$. In particular it produces no positive
Steenrod square. This does not classify arbitrary natural
transformations chosen only after passage to $D(k)$.

\emph{Proof.}
First let $X$ have finitely many nondegenerate simplices.
The module of simplicial maps into $A_\bullet^q$ is a finite
limit, using those simplices and their finitely many faces.
Exact additive functors preserve finite limits, so
\[
 L_q q\operatorname{Hom}(X,A^r)
       \simeq \operatorname{Hom}(X,L_q qA^r).
 \tag{B-12}
\]
On representable simplices, a strict natural transformation is
therefore a simplicial $k$-linear map of the reconstructed
underlying differential modules. Since there are only $n$ exterior
degrees on an $n$-simplex, all its components are well defined.

Before quotienting, $A^r=\bigoplus_{w\geq r}A_w^r$ has simplicial
homology only in weight $r$; this follows from the contractions in
B-05 and from $A_r^r=Z_r^r\simeq K(k,r)$. That homology is bounded
in weight, so $q$ kills it. Exactness of $L_q q$ shows that
$L_q qA^r$ is simplicially acyclic for every $r$.
The exact de Rham resolution and the identity $L_q qS=k$
therefore give
\[
 L_q qZ^i\simeq K(k,i)
\]
in simplicial homotopy over $k=\mathbf F_2$. Here all normalized
modules are vector spaces, so their acyclic summands split, exactly
as in B-05. Every component changing exterior degree is consequently
nullhomotopic. The reconstructed surjection from degree $i-1$
forms to degree $i$ cycles is a Kan fibration with contractible
source, so the lifting argument of B-06 again identifies the
cohomology of sections with homotopy classes into those cycles.
The degree-zero component fixes the unit at
$\Delta[0]$; naturality of dimension shifting then makes its action
the identity in every cohomological degree. By (B-12) the same is
true on every finite $X$.

For arbitrary $X$, the induced maps on ordinary cohomology are
natural transformations. The universal simplicial module
$K(\mathbf F_2,i)$ has finitely many simplices in each dimension,
by the explicit Dold--Kan formula. Its $(j+1)$-skeleton is a finite
simplicial set, and restriction to it detects a degree-$j$
cohomology class. Hence a positive-degree operation already known
to vanish on finite simplicial sets vanishes on the universal class,
and therefore on every class on every $X$. The same argument
identifies the degree-zero operation with the identity.
\hfill$\square$

\paragraph{B-11. Refuted in its exclusive reading; precise replacement for L8.}
The fact that $\Omega$ varies with the coefficient ring does not
turn its operations on constant-coefficient cohomology of spaces
into the nontrivial additive-group operations of scan 374.
The coefficient-fixed subgroup acts only by the identity, by
B-06, and the strict $\tau_\omega$ automorphisms of B-10 likewise
act by the identity over $\mathbf F_2$. This is the unit operation,
which belongs to both operation algebras when each is defined;
there is no positive-degree operation to place exclusively in one
of them. In particular the identity automorphism over $\mathbf F_p$
satisfies every automorphism hypothesis and induces the unit
Steenrod operation. It is a complete counterexample to the literal
assertion that the induced operation is excluded from the Steenrod
algebra. The distinction between the two coefficient problems
remains essential; it is not a mechanism producing an additional
operation.

For the full $\Omega_{\rm cont}(k)$, (B-10) is a semilinear action on
$\widehat S$-valued weighted cohomology. An action on ordinary
$k$-cohomology requires specifying how those coefficients are
removed. On the lowest-weight quotient the answer is
$x\mapsto\Delta^i x$ in degree $i$. For $k=\mathbf F_p$ and
$\Delta\ne1$ this family does not commute with suspension, so it
is not a single stable degree-zero Steenrod operation: suspension
from $\widetilde H^0(S^0;k)$ to $\widetilde H^1(S^1;k)$ would have
to intertwine multiplication by $1$ with multiplication by $\Delta$. It is
multiplicative as a degree-dependent rescaling. Over $\mathbf F_2$
that rescaling is the identity.

On the additive-group side the unit must not be confused with
$P^0$: Breen, \emph{Extensions du groupe additif}, \S4,
pp.~84--85, distinguishes the identity (the empty monomial) from
$P^0=F$, the Frobenius. His Theorem 4.6 and equation (4.13), p.~89,
give the Frobenius-semilinear coefficient relation. The operations
here fix $k$; they do not turn $a$ into $a^p$ on a general
$\mathbf F_p$-algebra. The fact that these two coefficient maps
agree on the field $\mathbf F_p$ itself does not identify the
natural transformations on all coefficient rings.

There is a precise comparison for the different algebra of
operations arising from higher equivariant products. May's
Definition~2.1, p.~160, specifies that extra structure; his
discussion on p.~154 explains why its $P^0$ is not generally the
identity \cite{BMay1970}. For the generalized operation algebra
$\mathcal B$ acting on $E_\infty$-algebras, Mandell's
Theorem~1.4, author PDF p.~5, proves
$\mathcal B/(1-P^0)\cong\mathcal A_{\rm Steenrod}$
\cite{BMandell2001}. Neither theorem identifies this
$\mathcal B$ with Breen's additive-group extension algebra, nor
turns these higher operations into the unary automorphisms
classified above. Our explicit computation B-06 supplies the
answer for those automorphisms.

\subsection{B.4. The rescaling monoid and its intersection with H}

\paragraph{B-12. Proved: the exact monoid, its units, and the
square-zero intersection.}
Let $S=\Gamma_k(T)$, with finite divided-power polynomials, and
let $G_k$ mean the monoid on scans 206--208. It is exactly the set
$S$ with law
\[
 (\lambda\star\mu)(T)=\lambda(T)\mu(T\lambda(T)),
 \qquad u_\lambda(T)=T\lambda(T),
 \tag{B-13}
\]
where substitution is divided-power substitution. Its identity is
$1$, and its group of units consists exactly of the ordinary ring
units $S^*$, with law (B-13). The continuous completed version has
underlying set $\widehat S$ and units exactly the series with
invertible constant coefficient.

If $H_k$ is the scalar-fixing endomorphism monoid, then
\begin{align*}
 G_k\cap H_k
  &=1+\operatorname{Ann}_S(T),\\
 \operatorname{Ann}_S(T)
  &=\bigoplus_{i\geq1}{}_{{i+1}}k\,T^{(i)},
       &\operatorname{Ann}_S(T)^2&=0.
 \tag{B-14}
\end{align*}
Thus the intersection is canonically this additive torsion group.
Replace the direct sum by the product for $\widehat S$. In
particular the intersection is trivial exactly when the additive
group of $k$ is torsion free.

\emph{Proof.}
The rescaling sends $U_r\mapsto\lambda(T)U_r$ and
$dU_r\mapsto\lambda(T)dU_r$; composition gives (B-13).
The parameter is detected by the coefficient of $U_0$, so the map
from this monoid to endomorphisms is injective.
A unit in the endomorphism monoid must induce an isomorphism on
the quotient by positive exterior degree and on its first exterior
piece. On the one-simplex the latter map is multiplication by
$\lambda$, after the coefficient-ring automorphism. Hence
$\lambda$ is a unit in $S\{X\}$ and evaluation at $X=0$ shows
$\lambda\in S^*$.

Conversely write $\lambda=a+\sum_{i>0}a_iT^{(i)}\in S^*$.
Then $a\in k^*$. The graded dilation homomorphism
$S\to S[z]$, $T^{(i)}\mapsto z^iT^{(i)}$, sends $\lambda$ to
a polynomial unit. Reducing modulo each prime ideal of $S$ shows
that every positive homogeneous coefficient $a_iT^{(i)}$ is
nilpotent. They generate a finitely generated nilpotent ideal $J$.
Remove the scalar dilation $T\mapsto aT$ from $u_\lambda$.
The resulting endomorphism $v$ is the identity modulo $J$, and
$v(a_iT^{(i)})-a_iT^{(i)}\in J^2$, since substitution sends a
homogeneous term to itself times a power of $1+J$. The product rule
then gives $(v-1)(J^r)\subset J^{r+1}$. Therefore $v-1$ is
nilpotent and the finite geometric sum is its inverse as a linear
map. A bijective ring map has a ring inverse, which is again a
PD substitution. Hence $u_\lambda$ is invertible. Its inverse
rescales $U_r$ by $u_\lambda^{-1}(\lambda)^{-1}$, proving that
$\lambda$ is a monoid unit. In the completion the same argument
uses the convergent weight filtration instead of a nilpotent ideal.

The coefficient equation $(\lambda-1)T=0$ is
$\lambda_0=1$ and $(i+1)\lambda_i=0$ for $i\geq1$, proving the
module formula in (B-14). We now verify its multiplication, which
is not supplied by just the parameter count on scan 208.
For $i,j\geq1$,
\[
 \gcd(i+1,j+1)\mid\binom{i+j}{i}.
 \tag{B-15}
\]
Indeed if $p^r$ divides both $i+1$ and $j+1$, then for every
$1\leq s\leq r$ the $s$th summand in Legendre's valuation formula
\[
 v_p\binom{i+j}{i}
 =\sum_{s\geq1}\left(
 \left\lfloor\frac{i+j}{p^s}\right\rfloor-
 \left\lfloor\frac{i}{p^s}\right\rfloor-
 \left\lfloor\frac{j}{p^s}\right\rfloor\right)
\]
is one, and all further summands are nonnegative. This proves
(B-15). If $(i+1)b_i=(j+1)c_j=0$, their product is annihilated by
$\gcd(i+1,j+1)$, by Bezout, and so
$\binom{i+j}{i}b_ic_j=0$. Every product of two elements in
$\operatorname{Ann}(T)$ is zero. In the completed case each
coefficient is still a finite sum, so the same conclusion holds.

For $\lambda\in1+\operatorname{Ann}(T)$ the scalar map $u_\lambda$
is the identity, since it fixes $T$ and its divided powers.
Thus (B-13) is ordinary multiplication, and (B-14) makes it
addition of the parameters. All these elements are units, with
inverse $1-b$ for $\lambda=1+b$. Finally nonzero torsion killed
by an integer $n\geq2$ appears in the $(n-1)$st summand, proving
the last assertion.\hfill$\square$

\paragraph{B-13. Refuted: the suggested semidirect monoid
decompositions; corrected split group extension.}
The coefficient map $G_k\to\operatorname{End}_{k,\rm PD}(S)$ has
image exactly the substitutions
\[
 f(T)=\sum_{n\geq1}f_nT^{(n)},\qquad f_n\in n k.
 \tag{B-16}
\]
Thus it need not be surjective even when $k$ is torsion free:
over $\mathbf Z$, $T+T^{(2)}$ is not in its image. Over a
$\mathbf Q$-algebra it is a bijection of endomorphism monoids,
but the proposed semidirect product of endomorphism monoids is
still false. Take $k=\mathbf Q$, scalar map $u(T)=0$, $A=X(T-X)$,
and $B_i=0$. B-01 gives a nonzero endomorphism. The only rescaling
with scalar map $u$ is $\lambda=0$, and composing that rescaling
on either side with a scalar-fixing endomorphism sends every
$U_r$ to zero. It cannot give this $A$. This verifies every
hypothesis of the counterexample, including the zero endpoint
values of $A$.

For \emph{automorphism groups} the correct universal statement is
instead the split exact sequence
\[
 1\longrightarrow H_k^{\rm aut}\longrightarrow
 \operatorname{Aut}_{k,\rm diff,simp}(A)
 \longrightarrow\operatorname{Aut}_{k,\rm PD}(S)
 \longrightarrow1,
 \tag{B-17}
\]
where the section sends $f(T)$ to the coordinatewise substitution
$U_r\mapsto f(U_r)$, with $B_i=0$. The same statement holds
continuously after completion. This section need not belong to
$G_k$. Consequently (B-17), rather than rescalings alone, gives
the semidirect product in all characteristics and with all torsion.

\emph{Proof.}
The identity $T T^{(n-1)}=nT^{(n)}$ proves (B-16).
The counterexample follows from B-01; $A\ne0$ is witnessed by its
coefficient of $XT$, while its values at $X=0,T$ are zero.
Any simplicial automorphism induces a scalar automorphism and,
since $S^+=(T)_{\rm PD}$, B-01 makes it a PD automorphism.
Conversely the coordinatewise substitution just specified is
simplicial and differential by B-01. If $g$ is the inverse PD
substitution, its coordinatewise version is the inverse.
Composition of these sections agrees with coefficient
composition. This proves split exactness, with kernel exactly
the scalar-fixing group.\hfill$\square$
